\appendix

\section{Proofs of the Schedulability and Worst-Case Guarantees}
\label{sec:proofs}

\begin{proof}[Proof of Theorem~\ref{thm:sched}]
\emph{Sufficiency} is constructive: allocate the standard floor $a_i = \beta b_i^{\mathrm{att}}$ to each $i\in\mathcal{S}$ and give each premium job its full requirement $a_i = b_i^{\mathrm{att}}$; by \eqref{eq:frontier} the total does not exceed $B$, and by \eqref{eq:attain} each premium job satisfies $T_i(a_i) \le c_i T_i^{\mathrm{solo}}$ even if all jobs communicate simultaneously.

\emph{Necessity} is worst-case: suppose $\sum_{\mathcal{P}} b_i^{\mathrm{att}} > B - \beta\sum_{\mathcal{S}} b_i^{\mathrm{att}}$. After honoring the standard floor, less than $\sum_{\mathcal{P}} b_i^{\mathrm{att}}$ remains, so some premium job $j$ receives $a_j < b_j^{\mathrm{att}}$. The periodic workload of \S\ref{sec:periodic-origin} admits the simultaneous-communication instant, at which $j$'s communication time stretches to $D_j/a_j > D_j/b_j^{\mathrm{att}}$ and its iteration time exceeds $c_j T_j^{\mathrm{solo}}$. Hence no scheduler can guarantee all premium SLOs in the worst case when \eqref{eq:frontier} fails.
\end{proof}

\begin{proof}[Proof of Theorem~\ref{thm:guarantee}]
The allocator reserves the standard floor first, leaving the premium pool $B_{\mathcal{P}} = B - \beta\sum_{\mathcal{S}} b_i^{\mathrm{att}}$, and distributes it by the equal-attainment waterline $a_i = \min(b_i^{\mathrm{att}}, \lambda b_i^{\mathrm{att}})$ with $\sum_{\mathcal{P}} a_i = B_{\mathcal{P}}$. Because all caps sit at the same relative level ($\lambda$), the waterline solves to \eqref{eq:lambda} directly. No premium job can be pushed below this level: pushing any job lower would free less than $\lambda b_i^{\mathrm{att}}$ and the waterline condition would be violated.
\end{proof}

\section{Staggering Relaxation: Detailed Treatment}
\label{sec:stagger-details}

This appendix provides the probabilistic derivation underlying Conjecture~1 in \S\ref{sec:guarantee}. The purpose is not to establish a formal theorem (which would require assumptions that do not hold under adversarial arrivals), but to quantify why the worst-case bound in Theorem~2 is systematically pessimistic in practice.

\subsection{Phase Model}

Let each job $i$ have a communication duty cycle $\phi_i = T_i^{\mathrm{comm,solo}} / T_i^{\mathrm{solo}} \in (0,1)$. Assume, for the purpose of this analysis, that phase offsets are independent and uniformly distributed over $[0, T_i^{\mathrm{solo}}]$. At a random instant, job $i$ is in its communication phase with probability $\phi_i$. The probability that all $|\mathcal{P}|$ premium jobs are communicating simultaneously is:
\begin{equation}
\Pr(\text{all premium jobs in communication phase}) = \prod_{i\in\mathcal{P}} \phi_i.
\end{equation}

For typical training jobs, $\phi_i$ ranges from $0.1$ to $0.3$ (see Table~\ref{tab:duty}). The probability of full simultaneity is therefore extremely small:
\begin{equation*}
\Pr \le (0.3)^{|\mathcal{P}|} = 10^{-|\mathcal{P}| \log_{10}(1/0.3)} \approx 10^{-0.52|\mathcal{P}|}.
\end{equation*}
For $|\mathcal{P}| = 8$ (the E1/E3 workloads), this is $\sim 10^{-4.2}$; for $|\mathcal{P}| = 10$, $\sim 10^{-5.2}$.

\begin{table}[h]
\centering
\caption{Duty cycles for common DNN models (measured on 2-node RoCEv2 testbed).}
\label{tab:duty}
\footnotesize
\begin{tabular}{@{}l c c c@{}}
\toprule
\textbf{Model} & $\phi_i$ & $T_i^{\mathrm{solo}}$ (ms) & $T_i^{\mathrm{comm,solo}}$ (ms) \\
\midrule
GPT-2 Tiny & 0.02 & 180 & 3.6 \\
GPT-2 Medium & 0.08 & 250 & 20 \\
ResNet-50 & 0.15 & 80 & 12 \\
LLaMA-2-7B & 0.25 & 320 & 80 \\
V6 AllReduce (2048MB) & 0.90 & 100 & 90 \\
\bottomrule
\end{tabular}
\end{table}

\subsection{Expected vs. Worst-Case Demand}

The expected simultaneous premium demand is:
\begin{equation}
\mathbb{E}\bigl[\text{simultaneous demand}\bigr] = \sum_{i\in\mathcal{P}} \phi_i b_i^{\mathrm{att}}.
\end{equation}
This is strictly less than $\sum_{i\in\mathcal{P}} b_i^{\mathrm{att}}$ for all $\phi_i < 1$. The difference is bandwidth that is reserved in the worst-case allocation but unused in expectation; work-conserving enforcement redistributes it to active jobs.

\subsection{Empirical Validation}

Figure~\ref{fig:e1-ladder} shows the gap between $\lambda$ (the worst-case waterline) and realized premium attainment across the E1 bandwidth sweep. Even when $\lambda$ falls below $0.5$, realized attainment remains above $0.8$; the gap grows with job count and shrinks as duty cycles increase. This confirms that the worst-case bound is conservative and that phase staggering is the primary mechanism driving the gap.


\section{Model Design Decisions}
\label{sec:design-notes}

\subsection{Why Not Model Computation--Communication Overlap Exactly?}

Exact overlap modeling would require runtime visibility into:
\begin{enumerate}
\item PyTorch DDP's gradient bucket partitioning (which depends on the model architecture and `bucket\_cap\_mb`),
\item NCCL's collective algorithm selection (ring vs. tree vs. collnet, determined by topology and message size),
\item The interleaving of communication kernels with backward pass kernels on the GPU, which depends on CUDA stream scheduling and is not exposed to the host,
\item Overlap efficiency, which varies across GPU generations and driver versions.
\end{enumerate}
Even if all this information were available, the schedulability test would become a mixed-integer non-linear program with no closed-form solution; recomputing it at every window boundary would be infeasible. The conservative serial model, by contrast, yields a linear test and an $O(N)$ allocator, and its safety is verified empirically across a range of overlap factors (Figure~\ref{fig:e11_overlap} in \S\ref{sec:eval}).

\subsection{Why Not Model Multiple Bottlenecks?}

The multi-link problem requires solving a multi-commodity flow (MCF) problem: each job's AllReduce traffic consumes bandwidth on multiple links simultaneously, and allocations on one link constrain allocations on another. The general MCF problem does not admit a closed-form feasibility test; solving it online at window boundaries would require iterative LP solvers, which are incompatible with the sub-millisecond recomputation target.

We choose single-link as the tractable first step for three reasons:
\begin{enumerate}
\item The single-link test is a \emph{component} that can be applied independently to each oversubscribed link in a multi-tier topology;
\item The spine uplink is empirically the dominant bottleneck in production clusters~\cite{Philly};
\item Satisfying the test on every link is a \emph{sufficient} condition for end-to-end SLO satisfaction (see \S\ref{sec:discussion}), so the test never admits an infeasible SLO set.
\end{enumerate}
Extending the frontier to coupled multi-link bottlenecks is an important direction for future work.

\section{Supplementary Evaluation Details}
\label{sec:app-supp}

\begin{table}[h]
\centering
\caption{24-job anchor at 400\,Gbps ($C^*>B$), ten seeds.}
\label{tab:anchor}
\footnotesize\begin{tabular}{lccc}\toprule
Policy & Mean SAS & SLO Attn (\%) & Collapse (\%) \\
\midrule
LongLiu & 0.8167 & 31.3 & 29.6 \\
LL-S & 0.8817 & 30.0 & 14.2 \\
SRPT & 0.8091 & 23.3 & 11.2 \\
Fair & 0.7639 & 21.7 & 1.7 \\
CRUX & 0.7590 & 21.2 & 5.8 \\
DF & 0.7723 & 20.4 & 1.2 \\
CASSINI & 0.4249 & 8.8 & 2.1 \\
\bottomrule\end{tabular}\end{table}

\begin{table}[h]
\centering
\caption{Lingjun 2023 trace replay results ($n{=}10$).}
\label{tab:trace-replay}
\footnotesize
\setlength{\tabcolsep}{3pt}
\begin{tabular}{@{}l c c c c@{}}
\toprule
Policy & P-attn (\%) & SAS (mean) & SAS (min) & $p$ vs.\ LongLiu \\
\midrule
LongLiu & $75.0\pm11.2$ & $1.179\pm0.083$ & $0.731\pm0.269$ & --- \\
LL-S & $85.0\pm18.3$ & $1.835\pm0.501$ & $0.523\pm0.609$ & $2.0{\times}10^{-1}$ \\
SRPT & $48.6\pm6.6$ & $1.217\pm0.168$ & $0.380\pm0.220$ & $3.8{\times}10^{-5}$ \\
DF & $40.0\pm19.1$ & $1.104\pm0.160$ & $0.541\pm0.222$ & $1.1{\times}10^{-3}$ \\
Fair & $35.7\pm7.5$ & $1.038\pm0.153$ & $0.370\pm0.148$ & $1.1{\times}10^{-6}$ \\
CRUX & $34.3\pm8.8$ & $1.003\pm0.155$ & $0.400\pm0.152$ & $1.1{\times}10^{-6}$ \\
CASSINI & $20.7\pm5.3$ & $0.815\pm0.099$ & $0.277\pm0.099$ & $9.9{\times}10^{-8}$ \\
\bottomrule
\end{tabular}
\end{table}

\subsection{CPU Emulator: Multi-Job Contention}
\label{sec:app-emulator}

We additionally validate on a CPU-based RDMA emulator running on the same two-node RoCEv2 testbed. Scenario S1 (5 jobs, 2P/3S) sweeps four regime depths; S2 (6 jobs, 3P/3S) tests starvation conditions. Each regime repeats for 5 rounds with alternating scheduler order.

\Name improves P-attn by $36\%$ over Static in the moderate regime ($0.28$ vs.\ $0.44$) and by $8\%$ in the deep regime ($0.57$ vs.\ $0.62$); all policies converge in the ample regime. In S2, \Name reduces J2's slowdown to $1.43\times$ against Static's $1.97\times$ ($28\%$ improvement), with bounded standard degradation ($\le 1.16\times$).

\begin{figure}[h]
\centering
\includegraphics[width=0.8\columnwidth]{figure/fig_expC_s1_pattn.pdf}
\caption{S1 P-attn across regimes.}
\label{fig:app-s1}
\end{figure}

\begin{figure}[h]
\centering
\includegraphics[width=0.8\columnwidth]{figure/fig_expC_s2_perjob.pdf}
\caption{S2 per-job slowdown.}
\label{fig:app-s2}
\end{figure}

\subsection{Robustness and Sensitivity}
\label{sec:app-robustness}

\paragraph{Scale of contention.}
We scale contending jobs from 14 to 100 (Figure~\ref{fig:app-scale}). Both \Name and DF reach $97.5\%$ at 14 jobs, but the gap widens with scale: at 50 jobs, \Name achieves $60.0\%$ vs.\ DF's $35.7\%$; at 100 jobs, $36.8\%$ vs.\ $6.0\%$. DF's scalar $\pi$ cannot encode per-job requirements when tens of jobs contend; \Name's closed-form allocator computes each job's rate before quantization, so the allocator's structure dominates.

\begin{figure}[h]
\centering
\includegraphics[width=0.8\columnwidth]{figure/fig_e12_dscp.pdf}
\caption{Scale of contention: \Name vs.\ DF.}
\label{fig:app-scale}
\end{figure}

\paragraph{Window size.}
$W\in\{5,10,20,50\}$ under dynamic workload (Figure~\ref{fig:app-window}). P-attn ranges from $69.1\%$ ($W{=}10$) to $76.4\%$ ($W{=}20$), within one standard deviation at every $W$. $W{=}20$ is the default; the experiment confirms the choice is not harmful.

\begin{figure}[h]
\centering
\includegraphics[width=0.8\columnwidth]{figure/fig_e13_window.pdf}
\caption{Window size $W$ sensitivity.}
\label{fig:app-window}
\end{figure}

\paragraph{Anchor freezing.}
Under sustained congestion, a control-loop baseline updating EMA only at highest priority exhibits hard freezing: 9--10 of 30 premium jobs frozen, $53.3\pm11.5\%$ P-attn. Active probing falls below baseline; passive recalibration narrows variance without statistically significant gain ($+2.0$~pp, $p=0.66$). \Name's closed-form allocator achieves $83.7\pm6.9\%$ with zero frozen jobs (Figure~\ref{fig:app-anchor}, Table~\ref{tab:app-anchor}).

\begin{figure}[h]
\centering
\includegraphics[width=\columnwidth]{figure/fig_e14_anchor.pdf}
\caption{Anchor freezing: remedies vs.\ \Name{} closed form.}
\label{fig:app-anchor}
\end{figure}

\begin{table}[h]
\centering
\caption{Anchor-freezing remedies.}
\label{tab:app-anchor}
\footnotesize
\begin{tabular}{@{}l c c c@{}}
\toprule
Config & P-attn (\%) & Hard-frozen & Seeds \\
\midrule
Baseline (control loop) & $53.3\pm11.5$ & 9--10 / 30 & 10 \\
Naive active probe & $48.3\pm2.4$ & 14 / 30 & 2 \\
Passive recalibration & $55.3\pm8.6$ & 16--17 / 30 & 10 \\
\Name (closed form) & $83.7\pm6.9$ & 0 / 30 & 10 \\
\bottomrule
\end{tabular}
\end{table}

\paragraph{Straggler resilience.}
We inflate compute times of a random job subset by factors $\in\{1.0,2.0,3.0,5.0\}$ at 400\,Gbps (Figure~\ref{fig:app-straggler}). \Name holds $72.5\%$ through $2\times$ and rises to $75.0\%$ at $5\times$; DF degrades from $52.5\%$ to $47.5\%$, and CRUX from $62.5\%$ to $52.5\%$.

\begin{figure}[h]
\centering
\includegraphics[width=\columnwidth]{figure/fig_e15_straggler.pdf}
\caption{Straggler resilience.}
\label{fig:app-straggler}
\end{figure}

\paragraph{Operator sensitivity to $\beta$.}
$\beta\in\{0.3,0.5,0.7,1.0\}$ at 500 and 630~Gbps (Figure~\ref{fig:app-beta}). Premium P-attn is essentially flat ($85$--$95\%$ at 500, $92.5$--$97.5\%$ at 630). Standard-tier S-cont improves monotonically from $0.978$ to $0.999$. Operators can choose $\beta$ freely within $[0.3,1.0]$ without compromising premium guarantees.

\begin{figure}[h]
\centering
\includegraphics[width=\columnwidth]{figure/fig_e16_beta.pdf}
\caption{Operator knob $\beta$ sensitivity.}
\label{fig:app-beta}
\end{figure}

\paragraph{Mixed-collective generality (E17).}
E17 replays the E1 14-job structure (8P/6S, FatTree $k{=}4$) with heterogeneous collectives: five AllReduce (DDP), four AllGather (ZeRO-3), three All-to-All (MoE), and two ReduceScatter (ZeRO-3) jobs ($n{=}10$, 600\,s). Per-job communication volumes are scaled to the collective ($0.5\times$ for the single-phase primitives). Table~\ref{tab:e17} reports P-attn across spine bandwidths. \Name leads at 400--630\,Gbps ($67.5\%$--$73.8\%$) and is statistically indistinguishable from the best baseline at 800\,Gbps where all policies converge; the E1 ordering is preserved, confirming that the guarantees are collective-agnostic.

\begin{table}[h]
\centering
\caption{Mixed-collective P-attn (\%), E17, mean$\pm$std over ten seeds.}
\label{tab:e17}
\footnotesize
\setlength{\tabcolsep}{4pt}
\begin{tabular}{@{}lcccc@{}}
\toprule
\textbf{Policy} & 400 Gbps & 500 Gbps & 630 Gbps & 800 Gbps \\
\midrule
LongLiu & $67.5\pm13.4$ & $71.2\pm15.6$ & $73.8\pm12.4$ & $73.8\pm12.4$ \\
LL-S & $46.2\pm26.4$ & $55.0\pm25.1$ & $60.0\pm24.2$ & $76.2\pm12.4$ \\
DF & $48.8\pm13.8$ & $55.0\pm18.8$ & $67.5\pm16.9$ & $77.5\pm19.4$ \\
SRPT & $48.8\pm21.6$ & $58.8\pm20.5$ & $70.0\pm16.9$ & $77.5\pm12.9$ \\
CRUX & $36.2\pm24.6$ & $41.2\pm21.3$ & $52.5\pm22.7$ & $70.0\pm17.9$ \\
CASSINI & $38.8\pm13.8$ & $45.0\pm12.1$ & $52.5\pm7.9$ & $60.0\pm9.9$ \\
Fair & $61.3\pm18.1$ & $65.0\pm19.4$ & $72.5\pm15.4$ & $77.5\pm11.5$ \\
\bottomrule
\end{tabular}
\end{table}

\paragraph{Component ablation (S3).}
We ablate the two components of \Name---the progress-deficit priority key and the closed-form $T$-target calibration---on the 24-job workload at 400\,Gbps ($n{=}10$, 600\,s). The full \Name{} reaches the highest SLO attainment ($35.0\pm9.9\%$, mean SAS $0.924$); freezing the $T$-target (LL-S) loses little ($32.9\%$, $p{=}0.13$), showing that the dynamic calibration mainly stabilizes variance rather than the mean. Replacing the deficit key with the static GPU-intensity key (CRUX, $24.6\%$) or with fair sharing ($25.8\%$) is clearly worse, and staggering communication phases---either CASSINI ($10.8\%$) or the full \Name{} with phase offsets ($12.9\%$)---collapses both attainment and throughput ($p<0.001$), because time-shifting dilutes the simultaneous-contention instants on which priority differentiation relies. The dynamic calibration is not free: it raises the handoff rate to $27.9\%$ (vs.\ $9.6\%$ for LL-S).

\section{Epoch-Boundary Control Flow}
\label{sec:app-epoch}

\begin{algorithm}[h]
\caption{Epoch-Boundary Control Flow}
\label{alg:epoch}
\begin{algorithmic}[1]
\State \textbf{Global:} $W$ (window size), $\alpha$ (EMA factor), communicators[7]
\Function{epoch\_start}{}
  \State $t_{\text{start}} \gets$ \texttt{time\_ns()}
\EndFunction
\Function{epoch\_end}{}
  \State $(D, T_{\text{comm}}, T_{\text{comp}}) \gets$ \texttt{ncclGetStats()}
  \State Update $T_i^{\text{comp}}$, $D_i$, $\mathrm{Bw}_{\text{ema}}$ per \S\ref{sec:anchor}
  \State $\{b_i^{\text{att}}\} \gets$ Equation~\eqref{eq:attain}
  \State $\{a_i\} \gets$ Priority from urgency $\pi_i$ (7-level monotonic mapping)
  \State $idx \gets$ \texttt{dscp\_to\_index}($a_i$)
  \State \texttt{select\_communicator}(communicators[$idx$])
\EndFunction
\end{algorithmic}
\end{algorithm}

\section{Additional Mechanism Figures}
\label{sec:app-figures}

\begin{figure}[h]
\centering
\includegraphics[width=0.9\columnwidth]{figure/fig5_pi_timeseries.pdf}
\caption{Per-job $\pi$ time series in E3$'$ (seed 0): (a)~pre-swap premium jobs, (b)~post-swap premium (red) and standard (gray) jobs.}
\label{fig:d1-transient}
\end{figure}

\section{Mechanism Boundary at the Wire}
\label{sec:app-boundary}

We examine when DSCP-based line-rate enforcement actually changes outcomes. A dual-flow raw-RDMA probe delivers strict-priority arbitration with a $253\times$ service ratio ($45.6$ vs.\ $0.18$\,Gbps). \texttt{ethtool -S} counters confirm \texttt{trafficClass} maps to distinct hardware queues. But with total demand ($18$--$37$\,Gbps) below the $50$\,Gbps wire, switch egress does not become the bottleneck; priority arbitration has little to arbitrate. \Name{} is faster on the tight job ($+3.5\%$ at 15\,Gbps, $+14.9\%$ at 30\,Gbps) but slower on the loose job ($-8.6\%$, $-3.9\%$). A GPT-2 training task shows the same boundary: the scheduler triggers P6 but communication does not improve because the bottleneck is host-side (NCCL proxy/PCIe). DSCP-based guarantees are necessary but not sufficient; they arbitrate only switch egress, and help precisely when the wire is the bottleneck.

\begin{figure}[h]
\centering
\includegraphics[width=0.9\columnwidth]{figure/fig7_v6_trajectory.pdf}
\caption{Mechanism boundary: window-by-window slowdown (mean over two rounds) for 15/30\,Gbps background $\times$ Job A/B (tight SLO).}
\label{fig:v6-boundary-traj}
\end{figure}
